\section{FPGA Architecture}
\label{sec:fpga_architecture}

The FPGA design is implemented in SystemVerilog and organized as a set of dedicated modules. The top-level file is \code{DE2\_115.sv}. The final design uses the FPGA for all real-time localization and control computations.

\subsection{Module Overview}

\begin{table}[H]
    \centering
    \caption{Important FPGA modules.}
    \label{tab:fpga_modules}
    \renewcommand{\arraystretch}{1.2}
    \small
    \begin{tabularx}{0.98\textwidth}{p{5.8cm} X}
        \toprule
        \textbf{Module} & \textbf{Function} \\
        \midrule
        \code{qmc5883p\_ctrl.sv} & Initializes and reads one QMC5883P magnetometer through I2C. \\
        \code{i2c\_master.sv} & Low-level I2C transaction controller. \\
        \code{mag\_calibration\_manager.sv} & Collects min/max samples and generates offset/scale coefficients. \\
        \code{mag\_calibrator.sv} & Applies offset and scale correction to raw magnetometer data. \\
        \code{carrier\_reference\_75hz.sv} & Generates sine/cosine reference samples for 75 Hz extraction. \\
        \code{carrier\_reference\_45hz.sv} & Generates sine/cosine reference samples for 45 Hz extraction. \\
        \code{mag\_lockin\_vector\_75hz.sv} & Performs rolling-window lock-in filtering and vector $H^2$ calculation. \\
        \code{h2\_average\_3sensor.sv} & Averages the latest 10 complete $H^2$ frames before classification. \\
        \code{h2\_lut\_classifier\_3sensor.sv} & Compares live normalized $H^2$ features with LUT ROM entries. \\
        \code{platform\_tracker\_controller.sv} & Converts local key decisions into one-key platform movement commands. \\
        \code{stepper\_drv8825\_controller.sv} & Generates STEP, DIR, and ENABLE\_N signals for the DRV8825. \\
        \code{mag\_uart\_streamer.sv} & Streams $H^2$ and final key/press frames to the PC. \\
        \bottomrule
    \end{tabularx}
\end{table}

\subsection{Sensor Readout}

Each QMC5883P controller checks the chip ID, configures the sensor for continuous measurement, and repeatedly reads the X, Y, and Z magnetic-field registers. In the current configuration, the sensor is set to continuous mode, 200 Hz output data rate, and $\pm 2$ G range. The FPGA receives three signed 16-bit raw values from each sensor.

The controller emits a one-clock \code{sample\_valid} pulse after each fresh sensor sample. Since the three I2C controllers do not necessarily finish in the same clock cycle, the top-level design collects valid pulses into a mask. A new lock-in sample is accepted only after all active sensors have produced a new sample.

In implementation, the three sensor controllers run in parallel instead of being multiplexed through one shared I2C controller. The top-level \code{DE2\_115.sv} module latches each sensor result when its valid pulse arrives, then clears the collection mask only after all three sensors have updated. This prevents the lock-in filter from mixing a new sample from one sensor with old samples from the other sensors.

\subsection{Magnetometer Calibration}

The calibration stage compensates for sensor offset and axis-scale mismatch. During calibration, the user rotates the complete sensor assembly. The FPGA records the maximum and minimum value for each active sensor axis:

\begin{equation}
    \mathrm{offset} = \frac{\mathrm{max} + \mathrm{min}}{2}.
\end{equation}

The measured radius of each axis is

\begin{equation}
    r_{\mathrm{axis}} = \frac{\mathrm{max} - \mathrm{min}}{2}.
\end{equation}

The scale coefficient is derived by comparing each axis radius to a target radius:

\begin{equation}
    \mathrm{scale}_{\mathrm{axis}} =
    \frac{r_{\mathrm{target}}}{r_{\mathrm{axis}}}.
\end{equation}

The calibrated output is then

\begin{equation}
    B_{\mathrm{cal}} = (B_{\mathrm{raw}} - \mathrm{offset}) \cdot \mathrm{scale}.
\end{equation}

This correction is implemented on FPGA. No Python calibration is needed during the final runtime.

\subsection{Streaming Control Style}

Most of the FPGA pipeline is controlled by one-clock valid pulses. A module latches its input when the upstream \code{valid} signal is asserted, performs its internal multi-cycle calculation if needed, then emits a new \code{valid} pulse for the next stage. This style is used for sensor readout, lock-in output, $H^2$ averaging, LUT classification, platform tracking, and UART streaming. It keeps the modules independent while still allowing slower blocks, such as division and ROM scanning, to take multiple clock cycles.

\placeholderfigure{Insert a diagram of the FPGA datapath: I2C readout, calibration, lock-in, H2 averaging, classifier, platform tracker, UART.}{FPGA datapath architecture.}{fig:fpga_datapath}{figures/fpga_datapath.png}
